% Part: lambda-calculus
% Chapter: syntax
% Section: beta

\documentclass[../../../include/open-logic-section]{subfiles}

\begin{document}

\olfileid{lam}{syn}{bet}
\olsection{$\beta$-લઘુકરણ}
\sourcecorrection{OLLAM-029}{સ્થિર અંગ્રેજી મૂળનો
અધ્યાય-ટિપ્પણી અને \texttt{\string\olfileid} આ
વિભાગને ભૂલથી પરિચય અધ્યાયમાં મૂકે છે. તે વાક્યરચના
અધ્યાયમાંથી આયાત થાય છે; તેથી અહીં અધ્યાય-ઓળખ
\texttt{syn} કરી છે.}

$(\lambd[m][(\lambd[y][y]) m])$ જોઈએ ત્યારે તેનો
$\lambd[m][m]$ સાથે સંબંધ હોવો જોઈએ એવું સહજ લાગે છે:
બીજું પદ પહેલાને ``સરળ'' કરવાથી મળતું હોય એમ જણાય છે.
\emph{$\beta$-લઘુકરણ} આ વિચારને ઔપચારિક સ્વરૂપ આપે છે.

\begin{defn}[$\beta$-સંકોચન, $\bredone$] \ollabel{defn:betacontr}
  \emph{$\beta$-સંકોચન} ($\bredone$) પદો પરનો સૌથી
  નાનો રચના-સુસંગત સંબંધ છે જે નીચેની શરત સંતોષે છે:
  \[
    (\lambd[x][N])Q \bredone \Subst{N}{Q}{x} 
  \]
  $P \bredone Q$ હોય તો $P$નું $Q$માં
  \emph{$\beta$-સંકોચન} થાય છે એમ કહીએ. $(\lambd[x][N])Q$
  સ્વરૂપનું પદ \emph{સંકોચ્ય પદ (રેડેક્સ)} કહેવાય છે.
\end{defn}

\begin{prob} \ollabel{prob:def}
  \olref[lam][syn][alp]{defn:aconvone}માં બદ્ધ ચલનું નામ
  બદલવા માટે આપી હતી તેવી $\beta$-સંકોચનની સમકક્ષ
  આગમનાત્મક વ્યાખ્યાઓ વિગતે લખો.
\end{prob}
  
\begin{defn}[$\beta$-લઘુકરણ, $\bred$] \ollabel{defn:betared}
  \emph{$\beta$-લઘુકરણ} ($\bred$) પદો પરનો
  $\bredone$ને સમાવતો સૌથી નાનો સ્વવાચક અને
  પરંપરિત સંબંધ છે. $P \bred Q$ હોય તો $P$નું
  $Q$માં $\beta$-લઘુકરણ થાય છે એમ કહીએ.
\end{defn}

સંદર્ભ સ્પષ્ટ હોય ત્યારે $\bredone$ને બદલે $\redone$
અને $\bred$ને બદલે $\red$ લખીશું.

અનૌપચારિક રીતે, $M \bred N$ ત્યારે અને માત્ર ત્યારે જ
જ્યારે $\beta$-સંકોચનનાં શૂન્ય અથવા વધુ પગલાંથી
$M$ને $N$માં બદલી શકાય.

\begin{defn}[$\beta$-સામાન્ય]
જે પદનું આગળ $\beta$-સંકોચન થઈ ન શકે તેને
\emph{$\beta$-સામાન્ય} કહે છે.
\end{defn}

$M \bred N$ અને $N$ $\beta$-સામાન્ય હોય તો $N$ને
$M$નું \emph{સામાન્ય સ્વરૂપ} કહીએ. કોઈ પદનું સામાન્ય
સ્વરૂપ અનન્ય છે કે નહીં એવો પ્રશ્ન થાય; જવાબ હા છે,
જેમ આગળ જોઈશું.

કેટલાંક ઉદાહરણો જોઈએ.
\begin{enumerate}
\item આપણને મળે છે:
\begin{align*}
(\lambd[x][xxy]) \lambd[z][z] & \redone (\lambd[z][z])(\lambd[z][z]) y \\
& \redone (\lambd[z][z]) y \\
& \redone y
\end{align*}
\item પદને ``સરળ'' કરવાથી તે વધુ જટિલ પણ બની શકે છે:
\begin{align*}
(\lambd[x][xxy])(\lambd[x][xxy]) & \redone (\lambd[x][xxy])(\lambd[x][xxy])y \\
& \redone (\lambd[x][xxy])(\lambd[x][xxy])yy \\
& \redone \dots
\end{align*}
\item સંકોચન પછી પદ યથાવત પણ રહી શકે છે:
\[
(\lambd[x][xx])(\lambd[x][xx]) \redone (\lambd[x][xx])(\lambd[x][xx])
\]
\item વળી, કેટલાંક પદોનું એક કરતાં વધુ રીતે લઘુકરણ
  થઈ શકે છે. દાખલા તરીકે,
\[
(\lambd[x][(\lambd[y][yx]) z]) v \redone (\lambd[y][yv]) z
\]
બહારના અનુપ્રયોગનું સંકોચન કરીને મળે છે; અને
\[
(\lambd[x][(\lambd[y][yx]) z]) v \redone (\lambd[x][zx]) v
\]
અંદરના અનુપ્રયોગનું સંકોચન કરીને મળે છે. જોકે આ
કિસ્સામાં બંને પદો આગળ જઈને એ જ પદ $zv$માં લઘુકૃત
થાય છે.
\end{enumerate}

છેલ્લા ઉદાહરણમાં મળેલું સમાન અંતિમ પરિણામ સંયોગ નથી;
તે ``ચર્ચ--રોસર ગુણધર્મ'' નામનો લૅમ્બડા કલનનો ઊંડો
અને મહત્વનો ગુણધર્મ દર્શાવે છે.

\begin{digress}
  સામાન્ય રીતે પદનું $\beta$-લઘુકરણ કરવાની એક કરતાં
  વધુ રીતો હોય છે. તેથી લઘુકરણની ઘણી વ્યૂહરચનાઓ
  ઘડાઈ છે; તેમાં સૌથી સામાન્ય \emph{સહજ વ્યૂહરચના} છે.
  આ વ્યૂહરચના હંમેશા સૌથી ડાબા સંકોચ્ય પદનું સંકોચન કરે
  છે; અહીં સંકોચ્ય પદનું સ્થાન આખા પદમાં તેની શરૂઆતથી
  નક્કી થાય છે. આ વ્યૂહરચનાનો ઉપયોગી ગુણધર્મ એ છે કે
  કોઈ વ્યૂહરચનાથી પદનું સામાન્ય સ્વરૂપ મળી શકે ત્યારે
  અને માત્ર ત્યારે જ સહજ વ્યૂહરચનાથી પણ મળી શકે.
  આગળ જુદું ન કહીએ ત્યાં સુધી સહજ વ્યૂહરચના વાપરીશું.
\end{digress}

\begin{defn}[$\beta$-સામ્ય, $\equal$]
  \emph{$\beta$-સામ્ય} ($\equal$) નીચેના નિયમો વડે
  આગમનથી વ્યાખ્યાયિત સંબંધ છે:
  \begin{enumerate}
  \item $M \equal M$.
  \item જો $M \equal N$, તો $N \equal M$.
  \item જો $M \equal N$, $N \equal O$, તો $M \equal O$.
  \item જો $M \equal N$, તો $PM \equal PN$.
  \item જો $M \equal N$, તો $MQ \equal NQ$.
  \item જો $M \equal N$, તો $\lambd[x][M] \equal \lambd[x][N]$.
  \item $(\lambd[x][N])Q \equal \Subst{N}{Q}{x}$.
  \end{enumerate}
\end{defn}

પહેલા ત્રણ નિયમો આ સંબંધને સામ્ય સંબંધ બનાવે છે;
પછીના ત્રણ તેને રચના-સુસંગત બનાવે છે; અને છેલ્લો નિયમ
તેમાં $\beta$-સંકોચનનો સમાવેશ કરે છે.

અનૌપચારિક રીતે, બે પદો $\beta$-સામ્યવાળાં ત્યારે અને
માત્ર ત્યારે જ કહેવાય જ્યારે એકને $\beta$-સંકોચનનાં
અથવા $\beta$-સંકોચનના ``વિપરીત'' પગલાંનાં શૂન્ય કે
વધુ પગલાંથી બીજામાં બદલી શકાય. વિપરીત $\beta$-સંકોચન
એ રીતે વ્યાખ્યાયિત છે કે $M$નું વિપરીત $\beta$-સંકોચન કરીને
$N$ મળે ત્યારે અને માત્ર ત્યારે જ $N$નું $\beta$-સંકોચન
કરીને $M$ મળે.

ઉપરના નિયમો ઉપરાંત આપણે આ સંબંધમાં વધુ નિયમો ઉમેરીશું.
વિસ્તારેલા સામ્ય સંબંધને $\equal[X]$થી દર્શાવીશું,
જ્યાં $X$ ઉમેરેલો નિયમ છે.

\end{document}
